\chapter{1986 Nobel Prize in Economics}
\textbf{Laureates: }James Buchanan (USA)

\section*{Reason for Award}
Theory political decision-making public choice

\section{What They Did}
James Buchanan founded public choice theory applying economic principles to political processes analyzed constitutional economics and government failure studied voting rules and collective action Developed framework treating politicians bureaucrats voters as self-interested actors operating within institutional constraints His work showed democratic processes produce inefficient outcomes special interest capture rent-seeking bureaucracy expansion Government failure analogous market failure requiring constitutional constraints limiting discretionary power Buchanan emphasized rules governing decision-making choosing rules before knowing outcomes behind veil of uncertainty Constitutional political economy examines fundamental rules constraining collective choice promoting liberty equity

\section{Impact on Economic Thinking}
Buchanan fundamentally altered understanding government behavior showing democratic institutions susceptible to same self-interested motivations economic agents Public choice theory revealed how politicians seek votes bureaus expand budgets special interests capture regulations for rent extraction rather than public benefit His work exposed limitations majority rule revealed paradoxes voting cyclical preferences revealed why democratic processes may not yield socially optimal outcomes Constitutional economics sought rules constrain discretionary power prevent tyranny majority minority oppression Buchanan s analysis showed institutions matter choices constrained by chosen rules Collective action problems explained why inefficient policies persist despite widespread opposition His work justified constitutional restraints term limits balanced expenditure requirements spending caps influencing state governments federal reform movements

\section{Modern Perspectives Today}
Public choice theory remains indispensable political economy political science government studies Analytical framework explains legislative bargaining campaign lobbying regulatory capture Rent-seeking behavior understood through Buchanan lens Constitutional politics informed debates about fiscal responsibility debt limits balanced budget amendments State-level reforms implementing term limits spending constraints reflect Buchanan-inspired principles His work continues guide scholars examining political institutions governance quality public sector efficiency Political economists apply public choice analysis democratizing transitions authoritarian regimes evaluating electoral systems political corruption institutional design Buchanan s legacy visible ongoing scrutiny government actions questioning public service motivations assessing institutional constraints effectiveness
